\section{Global equilibrium and a slack service obligation}\label{sec:global}

The local first-order conditions in Section~\ref{sec:equilibrium} are not sufficient for equilibrium. Without a lower service bound, a large retail-demand shift can drive service in the losing direction toward zero. The resulting demand--frequency feedback can create multiple fulfilled-expectations continuations and profitable branch-jump deviations. We do not claim equilibrium-selection-independent nonexistence when the floor is removed. Such demand--frequency vicious and virtuous cycles are known in transit systems \citep{baryosef2013}; we use the service floor only to obtain a globally well-defined continuation while preserving the local feedback at the candidate equilibrium.

\subsection{An exact global-equilibrium witness}

The existence proof uses the normalization
\begin{equation}
t=F=1,\qquad c=0,
\end{equation}
and the primitives
\begin{equation}
M=\frac23,\qquad q=\frac13,
\qquad x^*=\frac{23}{40},
\qquad
\frac{w}{Ft}=\frac{22797\sqrt{7599}}{7364680}.
\label{eq:witness}
\end{equation}
At this point the unconstrained service share is
\begin{equation}
s^*=\frac{\sqrt{149}}{\sqrt{149}+\sqrt{51}}\approx0.630895,
\end{equation}
so both minimum-service constraints are strictly slack. For $M=2/3$ and $q=1/3$, the $R$-side floor begins to bind at $x=2/3$; over the binding region the waiting-cost difference is constant and $G_{q,x}=2$.

On the slack region $x\leq2/3$,
\begin{equation}
g(x)=2+A H_x(x,2/3).
\end{equation}
The minimum of this slope on the relevant interval is attained at the floor boundary and equals
\begin{equation}
g(2/3)
=\frac{47133952-341955\sqrt{7599}}{23566976}
\approx0.735137>0.
\label{eq:global_monotonicity}
\end{equation}
Hence $G_q$ is strictly increasing over the whole consumer interval. Every unilateral price choice therefore induces a unique shopper/operator continuation, with corner market shares when the unique zero lies outside $[0,1]$.

\input{tables/witness_summary}

At the witness, the own-price second-order conditions are strict and the two best-response slopes have opposite signs, as reported in Table~\ref{tab:witness}. Global optimality is established by reparameterizing each retailer's unilateral price problem by the uniquely induced share $x$. For $L$, the only competing maximum outside $x^*$ is the service-floor boundary, and the exact profit gap is
\begin{equation}
\pi_L(x^*)-\pi_L(2/3)
=\frac{358680877-4103460\sqrt{7599}}{1325642400}
\approx0.000734>0.
\label{eq:L_gap}
\end{equation}
For $R$, the best deviation in the binding-floor region is the vertex of a concave quadratic and remains below the candidate profit by approximately $0.009698$. Exact root isolation also verifies that no interior stationary point dominates $x^*$.

\begin{proposition}[Global network-mediated strategic asymmetry]\label{prop:global_asymmetry}
There exists a nonempty open set of primitives for which the model has a global pure-strategy retail-price Nash equilibrium satisfying
\begin{equation}
\frac{dBR_L}{dp_R}<0<\frac{dBR_R}{dp_L}.
\end{equation}
\end{proposition}

\noindent
The proof is constructive. Equations \eqref{eq:witness}--\eqref{eq:L_gap} provide an exact witness with a globally single-valued continuation, strict own-price second-order conditions, strict opposite reaction signs, and strict global-deviation profit gaps for both firms. All decisive inequalities are strict, so continuity yields a nonempty open neighborhood with the same equilibrium structure. A deterministic audit of a nearby parameter grid finds 595 candidates satisfying the local conditions, of which 444 also pass the monotone-continuation and global-deviation checks; this audit supports but does not replace the exact argument.

\subsection{Why a nonbinding service floor can matter globally}

The service obligation is not the source of Proposition~\ref{prop:global_asymmetry}. At the equilibrium it is slack, so neither the local demand slope nor the local price first-order conditions depend on $q$. Its role is off path: sufficiently large price deviations can otherwise trigger extreme service withdrawal from the rival direction.

\begin{proposition}[Slack service obligation and global support]\label{prop:slack_floor}
A minimum service obligation can be strictly slack at a price equilibrium satisfying Proposition~\ref{prop:global_asymmetry} while supporting that equilibrium by truncating only sufficiently extreme off-equilibrium service reallocations.
\end{proposition}

For the exact witness, the same on-path prices, market share, frequencies, and welfare are supported over the intermediate interval
\begin{equation}
q_L\approx0.324091<q<q_U\approx0.344228,
\label{eq:support_band}
\end{equation}
while the floor remains slack at the candidate. For $q$ just below $q_L$, while the continuation remains monotone, $L$ can cut price enough to exploit the high-$L$ continuation before the floor activates. If $q>q_U$ while still slack at the candidate, the off-path protection of $R$ becomes strong enough to create a profitable $R$ deviation. Thus a regulation that is nonbinding in normal operation can be too weak or too strong solely through its effect on off-equilibrium continuations.

Equation \eqref{eq:support_band} is witness-specific. We do not interpret $q_L$ and $q_U$ as universal policy thresholds, nor do we solve for an optimal service standard. Minimum service itself is a standard transit-service primitive; the result here concerns its nonbinding-on-path role in the downstream global price game.
